% Source: sources/returns/S6_04_MAPS_MANIFOLDS/sections/fourfolds_D2_critical_frame.tex:2–413
\paragraph{Theorem D.2.44 (Proved: the seven semistable residual-four and residual-eight smooth types).}
Consider the default semistable rows OS55 with $(a,b)=(5,4)$ and zero
index $2,3,4$, and OS45 with $(a,b)=(1,8)$ and zero index $1,2,3,4$.
Use the primitive Mumford cap at every multiplicative site, no additional
regular logarithmic transforms, and either simply connected ordinary
clutch $e\gamma$, $e=\pm1$. Write $c=b$.
The resulting oriented six-manifolds have
\[
 H_2(X;\mathbb Z)=H_3(X;\mathbb Z)=\mathbb Z/c,
 \qquad \lambda(X)=0,\qquad \mu(X)=0,
 \qquad
 P(X)=\begin{cases}2,&\mathrm{OS55},\\0,&\mathrm{OS45}.
 \end{cases}                                                   \tag{D44.1}
\]
Here $P$ is Zhubr's refined quarter-Pontryagin class, transported to
$\mathbb Z/c$ by any generator; both displayed values are independent
of that choice. In particular all three OS55 manifolds are oriented
diffeomorphic to $Z_{4,1}$, and all four OS45 manifolds are oriented
diffeomorphic to $Z_{8,0}$, with the following concrete surgery models.
The manifold $Z_{c,0}$ is surgery on an embedded two-sphere representing
$c$ times the generator of $H_2(S^2\times S^4)$.
The manifold $Z_{c,1}$ is surgery in
$M_{\rm conic}\#(S^2\times S^4)$ on an embedded two-sphere representing
$(t,-c/2)\in\mathbb Z/2\oplus\mathbb Z$, where $t$ is the generator
of the conic-surgery manifold of D.2.8. These are spin surgeries with
their unique normal framings. This theorem makes no assertion about
the remaining additive profiles or additional regular logarithms.

\emph{Proof: previously established invariants and the last binary choice.}
The marked homology calculation D.2.26 and its all-weight verification
D.2.41 give the stated integral groups, with no other reduced homology
except $H_6=\mathbb Z$. The full integral support-function and
semicharacter certificate
\path{code/fourfolds_D2_2power_AH.py}
verifies the integral line-descent and actual-angle hypotheses
used in D.2.35 for these rows. Here are the complete remaining
localization steps. The subgroup $G=C_c$ of the scalar circle has
fixed locus the rational curve $C$ and the one exceptional
elliptic double curve $E$ at $I_c$; it acts freely off them.
The actual normal characters are $(1,-1)$ on $C$ and $(w,-w)$
on $E$, with $w^2=1$ in $\mathbb Z/c$. The two normal lines on
$C$ have degree $-1$, and those on $E$ have degree zero.
The symmetric auxiliary line of D.2.35 has character $c/2$ on
$C$, zero on $E$, and degree coprime to $c$ on $C$.
Its identification with the global mod-$c$ class is the same
integral common-piece argument as there: the boundary-to-cap
surjections and primitive degree-one map make $H_2(W)\to H_2(X)$
onto; the line has zero swept values and kills the ordinary
clutch relation modulo $c$. Moreover $\pi_1(W_G)$ has the
unique torsion subgroup $G$. The difference between equivariant
lifts therefore has the same point character on $C$ and $E$.
After one base-character adjustment, the stated auxiliary
characters are the characters of a global equivariant lift.

Normalize its ordinary class $\alpha\in H^2(X;\mathbb Z/c)$
to have degree one on $C$. Put $t\in H^2(BG;\mathbb Z/c)$,
$x=[C]^*$ and $y=[E]^*$. Its restrictions have the form
\[
 \alpha|_C=x+(c/2)t,\qquad
 \alpha|_E=A y+\eta\zeta,\qquad
 e_G(N_C)=-t^2,\qquad e_G(N_E)=-w^2t^2,                       \tag{D44.0a}
\]
where $\eta$ has degree one on $BG$ and $\zeta$ lifts from
integral degree-one cohomology of $E$. Localization over
$\mathbb Z/c$ is legitimate: the complement is free, so its
cohomology vanishes after inverting $t$. The degree-one
localization identity gives $A=-w^2$. The elliptic restriction
has cube zero, because $y\zeta=0$ and $\zeta^2=0$; the rational
curve's contribution to the cube is $-3(c/2)^2=0\pmod c$.
Thus $\langle\alpha^3,[X]\rangle=0$ in the full coefficient
ring, not merely modulo two.
The equivariant complex determinant is trivial: its ordinary
class is zero and its total character on $C$ is zero.
Hence $\lambda_G=-c_{2,G}$. On either fixed curve,
$c_{2,G}(TX)=e_G(N)$, since the square of the tangent-line
class vanishes. Therefore
\[
 \langle\alpha\lambda,[X]\rangle=-(1+A)=-(1-w^2)=0.           \tag{D44.0b}
\]
The normalized class detects the cyclic degree-four group, so
$\lambda=0$. Zhubr's Sections 3.11--3.14 identify the complete
group $H_6(K(\mathbb Z/c,2);\mathbb Z)=\mathbb Z/c$ by the cube,
giving $\mu=0$ \cite[Sections 3.11--3.14]{manZhubr75}.
The free-coefficient calculation
and coefficient surjectivity in Zhubr's Sections 3.11--3.14 give
$2P=\operatorname{PD}(\lambda)$
\cite[Theorem 5 and Sections 3.11--3.14]{manZhubr75}.
Thus the two remaining possibilities
are exactly $P=0,c/2$. The calculation below distinguishes them by
the actual tangent framing, without halving a congruence or using
the conjectural change-of-three-frame formula of \cite{ffZhubr13}.

\emph{A framing parity that depends only on the three-cycle.}
We first record the precise primary construction used to compare a
regular three-torus with a torsion three-sphere.
Let a closed oriented six-manifold $Y$ have a stable tangent framing.
For an embedded oriented three-manifold $N\subset Y$ with a chosen
normal framing, denote by $[N]_{\rm fr}\in\Omega^{\rm fr}_3
\cong\mathbb Z/24$ its induced stable framed-bordism class. Then
\[
 [N]_{\rm fr}\pmod2
   =q_Y\bigl(\operatorname{PD}_{\mathbb F_2}[N]\bigr)             \tag{D44.2}
\]
for a function of the indicated mod-two cohomology class alone.
This statement includes torsion classes in integral homology.
Here is a derivation directly from Brown's construction
\cite[pp.369--372]{ffBrown72}. Set $K_3=K(\mathbb F_2,3)$ and let
$j:S^3\to K_3$ represent its fundamental class. Brown's $F(u,0)$
is literally the stabilization of the map $u:Y\to K_3$.
The framed Pontryagin--Thom collapse $Y\to S^3$, followed by $j$,
therefore represents $F(\operatorname{PD}[N],0)$.
Apply the framed duality map $W_*A_\alpha$ of his Theorem 1.7 and
Corollary 1.11. Its value is the image of $[N]_{\rm fr}$ under
\[
 j_*:\pi_3^{\rm S}\longrightarrow\pi_6^{\rm S}K_3
          \cong\mathbb Z/2.                                  \tag{D44.3}
\]
The target is Brown's Lemma 1.2. The remark immediately following
his Corollary 1.13 identifies $j_*\nu$ with its nonzero generator:
the coefficient is the mod-two Hopf invariant, and the quaternionic
Hopf element $\nu$ has Hopf invariant one. Dividing Brown's
$\mathbb Z/4$-valued function, whose values here are $0,2$, by two
gives (D44.2). This is a Pontryagin--Thom argument about the specified
normal framing; no compression of an arbitrary framed torus to a
sphere is being assumed.

For clarity we also fix the sphere convention. The homomorphism
$\pi_3SO(3)\to\pi_3SO(6)$ has image $2\mathbb Z$.
Consequently changing the normal framing of an embedded $S^3$
does not change the parity of its induced stable framed class.
For the standard sphere bounding a ball, the induced stable framed
class is zero and its tangent three-frame represents the nonzero
class in $\pi_3V_{6,3}=\mathbb Z/2$ relative to a constant ambient
three-frame. Equivalently, the tangent Lie framing on $S^3$ has
class $\nu$, and a change of ambient frame by degree $d$ changes
that class by $d\nu$, while the relative three-frame class changes
by $d\bmod2$. Thus Zhubr's required nontrivial relative
three-frame is equivalent to \emph{even} induced stable framed
class on the sphere. These conventions agree with the quaternionic
calculation in \cite[Section 5]{ffZhubr13}.

Since $\lambda(X)=0$, $TX$ is stably trivial: the only obstruction
for a spin bundle on a six-dimensional complex is its class in
$H^4(-;\pi_3\mathrm{Spin})$, namely $\lambda$.
Moreover $TX$ is actually trivial on $X^\circ=X\setminus\{p\}$.
For example its complex bundle has $c_1=c_2=0$ and its last
obstruction lies in $H^6(X^\circ;\mathbb Z)=0$.
Any framing there extends stably across $p$, since
$\pi_5SO=0$. All such full framings have the same relative
three-frame parity on an embedded torsion $S^3$: their difference
lifts to $\mathrm{Spin}(6)$ and its degree-three class is a class
of $H^3(X^\circ;\mathbb Z)$, which evaluates to zero on a torsion
three-cycle.
An embedded sphere in the integral class of the regular torus
$T=\langle e_1,e_2,\gamma\rangle$ used below can also be
constructed explicitly, so its existence is not an additional
representability assumption. Apply Kaplan's construction to
the spin structure induced on the normally framed $T$:
it gives a four-dimensional filling with one zero-handle and
only two-handles, preserving that almost framing
\cite[Definition 1.7, Remark 1.8, Theorem 3.1]{mapsKaplan79}.
Delete the zero-handle and reverse the cobordism. Realize the
resulting two-handle attachments as ambient surgeries in
$X^\circ$. Their attaching circles bound embedded disks
with interiors disjoint from $T$ and from one another, since
$\pi_1(X^\circ)=0$, $2+3<6$ and $2+2<6$.
Prescribe a boundary collar along the first normal vector.
The two tangent framing directions of each attachment and
the other two normal directions give a four-frame around
the boundary of its disk. Its extension obstruction in
$\pi_1SO(4)=\mathbb Z/2$ is exactly the spin obstruction
of that two-handle, and vanishes. Thickening these disks
therefore performs the prescribed surgeries inside $X^\circ$.
The resulting embedded $S^3$ has the same integral homology
class as $T$. Brown's identity already supplies its parity;
this construction additionally verifies its geometric
realization without changing the prescribed spin structure.
Zhubr's Theorem 1 now implies
\[
 P(X)=0\quad\Longleftrightarrow\quad
 \text{a full ambient framing is admissible on the torsion }S^3.
                                                               \tag{D44.4}
\]
Indeed the forward implication is not an assumption about arbitrary
three-frames. Take the admissible three-frame supplied by that
theorem. If $P=0$, its spin rank-three complement has vanishing
primary obstruction $c_2$; the only possible further obstruction
on the punctured six-manifold is in $H^5(X^\circ;\mathbb Z/2)=0$.
The complement is therefore trivial, and the three-frame extends
to a full framing. Conversely an admissible full framing has
trivial complement and hence refined class zero. Independence
of the full framing proves (D44.4) for the particular framing
used below.

\emph{The reference frame and the actual critical-cycle defect.}
Write $f:X\to\mathbb P^1$ for the torus fibration and
$\mathcal V=\ker(df)$ for its holomorphic differential kernel.
This is a locally free rank-two sheaf, including at the critical
curves. In the nodal chart $f=xy$ it has basis
$x\partial_x-y\partial_y,\partial_z$; in the triangular cusp
chart $f=xyz$ it has basis
$x\partial_x-y\partial_y,y\partial_y-z\partial_z$.
The scalar $\delta$-action defines a nowhere-zero section of this
\emph{abstract} vector bundle, although its image as a vector
field vanishes on the fixed curve. In a toric chart its two
coefficients are those of the nonzero primitive scalar lattice
vector, so they cannot both vanish. This description is unchanged
by the translation clutch.
As $H^2(X;\mathbb Z)=0$, the quotient line $\mathcal V/\mathcal O\delta$
is topologically trivial; hence $\mathcal V$ has a global smooth
frame whose first vector is $\delta$. The line
$\mathcal L=f^*T\mathbb P^1$ is also topologically trivial.

Let $Z$ denote the scheme defined by the coefficients of $df$.
There is an exact sequence
\[
 0\longrightarrow\mathcal V\longrightarrow TX
  \xrightarrow{df}\mathcal L
  \longrightarrow\mathcal L\otimes\mathcal O_Z
  \longrightarrow0.                                         \tag{D44.5}
\]
Its codimension-two cycle is the literal sum $C_{\rm crit}$ of
the complex double curves, each with coefficient one. This
includes the three axes at a triple point: the ideal
$(xy,xz,yz)$ has multiplicity one on each axis. Since
$\mathcal V$ and $\mathcal L$ are topologically trivial,
(D44.5) gives $c_2(TX)=[C_{\rm crit}]$. Thus this cycle is
integrally nullhomologous.

Off $Z$, splitting (D44.5) identifies $TX$ with the framed
bundle $\mathcal V\oplus\mathcal L$. Its comparison with a
global frame on $X^\circ$ has an integral degree-three class
$a$ whose local residue on every smooth double curve is one,
up to the same overall sign. In the model $xy$ this is the
explicit Bott matrix
\[
 \frac1{\sqrt{|x|^2+|y|^2}}
 \begin{pmatrix}x&\overline y\\-y&\overline x\end{pmatrix}
       \in SU(2).                                            \tag{D44.6}
\]
It has degree one on the normal three-sphere; adding the
tangent line to the critical curve sends this to the stable
generator, not twice that generator. The latter factor two
belongs to the different map $SO(3)\to SO$, as is also seen
from $\lambda=-c_2$ for this complex rank-two normal calculation.
Removing the finitely many triple points does not affect this
degree-three residue calculation. The integral relative
cohomology sequence, or equivalently the Thom intersection
construction, therefore gives
\[
 \langle a,[T]\rangle
      =\pm\operatorname{lk}(T,C_{\rm crit})                   \tag{D44.7}
\]
for any regular three-torus $T$ disjoint from the critical
cycle. The ambiguity of adding a class from $H^3(X^\circ;\mathbb Z)$
evaluates to zero on our torsion three-cycle. Thus the integer
in (D44.7), not merely a residue modulo $c$, is well defined.

Take specifically $T=\langle e_1,e_2,\gamma\rangle$ in a regular
fibre inside the ordinary clutch disc. The reference frame
$\mathcal V\oplus\mathcal L$ has zero degree-three difference
from the flat translation frame there. Indeed the vertical
frames have the same first vector $\delta$; their change lies
in the stabilizer $\mathbb C\rtimes\mathbb C^*$, which retracts
to a circle. The base line contributes another $\mathbb C^*$,
and changing the horizontal splitting contributes a contractible
space. None contributes a degree-three class. The induced flat
framing of $T^3$ is its product Lie framing, with stable class
$\eta^3=12\nu$, an even element of $\mathbb Z/24$.
A possible determinant-circle transition cannot change this
parity. A map $T^3\to S^1$ is homotopic to a character and,
after an integral change of torus basis, depends on just one
circle. Its real rotation can be placed stably in that circle
and one normal direction. It changes just that circle's framed
class, so the resulting product-framed torus has class either
$0$ or $\eta^3$, both even. After splitting off the determinant,
the remaining $SU(3)$ transition is a degree-$k$ map on a
three-dimensional complex and changes the framed class by
$\pm k\nu$. Thus its exact parity is the integer in (D44.7).
Equations (D44.2) and (D44.7) consequently give
\[
 q_X(\operatorname{PD}[T])
       =\operatorname{lk}(T,C_{\rm crit})\pmod2.              \tag{D44.8}
\]
The torus represents an odd generator of $H_3(X)=\mathbb Z/c$
by the explicit linking calculation below. Its reduction is
therefore the same as that of an embedded torsion generator
$S^3$. Combine (D44.4) and (D44.8) to obtain the actual decision
\[
 P(X)=\frac c2
       \bigl(\operatorname{lk}(T,C_{\rm crit})\bmod2\bigr).
                                                               \tag{D44.9}
\]

\emph{The integral linking calculation.}
All local preimages and bounding coefficients here are those
of the actual marked caps, not an invariant-form calculation
on an unmarked torus. At each site choose the recorded
determinant-one elliptic basis $(u_1,u_2)$ and use the short
and long planes $(u_1,\delta)$ and $(u_2,\gamma)$. The exact
valuation matrix is
\[
 L=\begin{pmatrix}n&-bH\\aH&R\end{pmatrix}.                  \tag{D44.10}
\]
At a rank-one site write $L=\alpha r$ with $\alpha$ primitive.
Choose a short complement $\sigma$ with
$\det(\alpha,\sigma)=1$ and long vectors $(\beta,\ell)$ with
$r\beta=\rho>0$, $r\ell=0$, and $\det(\beta,\ell)=1$.
The literal transverse complex elliptic curve has regular-fibre
preimage $\sigma\wedge\ell$. The product nodal chart supplies
the specialization chain with degree one. These signs follow
by orienting the vanishing elliptic factor as $(\alpha,\beta)$
and the transverse factor as $(\sigma,\ell)$.

At the unique rank-two cusp $\det L=1$. Put
$(b_1,b_2)=(u_2,\gamma)L^{-1}$ and $(a_1,a_2)=(u_1,\delta)$,
so $Nb_j=a_j$. The three integral support functions
$r(r-1)/2,rs,s(s-1)/2$ have degree matrix
\[
 \begin{pmatrix}0&-1&0\\1&1&-1\\-1&0&0\end{pmatrix}          \tag{D44.11}
\]
on the three complex opposite-pair double curves in cyclic
ray order $(1,0),(0,1),(-1,1)$. This unimodular matrix gives
degree vector $(-1,1,-1)$ for their sum. The fibre vector
\[
 a_1\wedge b_1+b_1\wedge a_2+a_2\wedge b_2                  \tag{D44.12}
\]
has precisely those degrees and zero long--long component.
The latter is the degree-two class of the dual torus, so these
four coordinates identify its specialization with the actual
sum of the three double curves integrally. In particular no
undetected fourth homology component is being suppressed.

Transport the five preimages to the common regular fibre along
the fixed base paths. Their sum is a representative of
$C_{\rm crit}$ in the exact marked presentation of D.2.41.
All transport and cap relations can be realized by chains
outside the ordinary disc. In that disc the four normal
Thom chains are the products of a lattice loop with the
disc, and the redundant fifth chain has loop vector
$(0,0,0,-e)$. The intersection of a bounding relation with
$T=e_1e_2\gamma$ is consequently its clutch $\delta$
coefficient. Integral Smith solutions give
\[
 \begin{array}{c|c|c|c}
  \mathrm{OS}&c&\text{zero indices}&
       \operatorname{lk}(T,C_{\rm crit})\text{ for }e=1\\\hline
  45&8&1,2,3,4&4,4,4,4\\
  55&4&2,3,4&37,7,27.
 \end{array}                                                  \tag{D44.13}
\]
For example, in OS55 at zero index $2$, the five cap vectors
in ordered bivector coordinates
$(e_1e_2,e_1\delta,e_1\gamma,e_2\delta,e_2\gamma,\delta\gamma)$
are
\[
 \begin{gathered}
 (0,0,0,-12,0,5),\quad(3,0,1,-12,0,4),\\
 (-19,-23,-10,0,5,6),\quad(0,0,0,0,0,1),\quad(0,4,0,0,0,1).
 \end{gathered}                                               \tag{D44.14}
\]
Their sum has an integral bounding relation with clutch loop
$(0,0,37,0)$. The same Smith calculation for $c$ times a
two-dimensional torsion generator gives a clutch coefficient
coprime to $c$ in every row. Nondegeneracy of the torsion
linking pairing proves that $T$ is an odd generator of
$H_3(X)$. Every integral relation among presentation columns
has zero clutch $\delta$ coefficient, so the displayed
linking does not depend on a selected Smith solution.
Changing the stated short complements or the cusp mixed
preimage by an actual cap-kernel class also changes that
coefficient by zero. For $e=-1$ the exact same calculation
with its actual clutch vector reverses the displayed integer
signs and retains the unit torsion-linking test.

The executable
\path{code/manifolds_even_critical_cycles.py}
retains all fourteen rows, the full integral coefficient
vectors, local bases, alternative preimages and both clutch
signs in
\path{records/manifolds_even_critical_cycles.json}.
Its integral computation is independent of the differential
and framing arguments above. Formula (D44.9) and table
(D44.13) prove (D44.1).
The separate executable
\path{code/fourfolds_D2_frame_certificate.py}
checks the local kernel and Bott-matrix identities, the
unimodular cusp degree matrix and, for all fourteen marked
rows, the independent degree identity
$\Omega(C_{\rm crit})=-be\,\operatorname{lk}(T,C_{\rm crit})$
for the actual pure fibre functional
$\Omega=a e_1^*e_2^*+b\gamma^*\delta^*$.
It retains its inputs' hash and every local degree in
\path{records/fourfolds_D2_frame_certificate.json}.

\emph{Recognition of the surgery models.}
The normal rank-four bundle of any embedded sphere in these
spin six-manifolds is trivial, and its framing is unique
since $\pi_2SO(4)=0$. The stated surgeries have $H_2=H_3=\mathbb Z/c$
and zero free third Betti number. The classifying map to
$K(\mathbb Z/c,2)$ extends across each surgery trace.
The model $Z_{c,0}$ has refined class and fundamental-class
invariant both zero. For $Z_{c,1}$, the conic summand's refined
class maps to $c/2$, and the complete cubic maps to zero for
$c=4,8$, since $(c/2)^3=0\pmod c$. Thus their full Zhubr
invariants are respectively $(P,\mu)=(0,0)$ and $(c/2,0)$,
using naturality in \cite[Theorem 5]{manZhubr75} and the full
conic invariant established in D.2.8.
The classification theorem \cite[Theorem 3]{manZhubr75}, including
its finite invariants, gives the asserted oriented diffeomorphisms. No conclusion
here rests on ordinary homology alone.\hfill$\square$
